\documentclass{article}
\usepackage{amsmath}
\usepackage{amsthm}
\usepackage{amsfonts}
\usepackage{amssymb}
\usepackage{bbold}
\usepackage{dsfont}
\usepackage{graphicx}
\usepackage{cleveref}
\usepackage{bm}
\usepackage{booktabs}
\usepackage{siunitx}
\usepackage{xcolor}
\usepackage{pdfpages}
\usepackage[utf8]{inputenc}
\usepackage[english]{babel}
\newtheorem{definition}{Definition}
\newtheorem{example}{Example}
\newcommand{\exP}[1]{\mathrm{e}^{#1}}
\newcommand{\dif}{\mathrm{d}}
\renewcommand{\vec}[1]{\mathbf{#1}}
\newcommand{\pdagger}{{\phantom{\dagger}}}
\newcommand{\state}[1]{\vert #1\rangle}
\newcommand{\cstate}[1]{\langle #1\vert}
\newcommand{\tr}[1]{\mathrm{tr}\left\{#1\right\}}
\newcommand{\trs}[2]{\, \mathrm{tr}_{#1}\{#2\}}
\newcommand{\expc}[3]{\langle #1\vert #2 \vert #3\rangle}
\usepackage[doi=false,url=false,isbn=false,backend=bibtex,style=nature,natbib=true,sorting=none]{biblatex}
\bibliography{bib.bib}
\usepackage{geometry}
\geometry{
	a4paper,
	total={170mm,257mm},
	left=15mm,
	right=15mm,
	top=15mm,
}
\begin{document}
\title{Notes}
\author{Rudolf Smorka}
\maketitle
\section{Relation between my model to the Dicke model}
\begin{itemize}
	\item Dicke model: One cavity mode interacts collectively with $N$ two-level systems
	\begin{equation*}
		H = \hbar \omega a^\dagger a + \hbar \omega_0 J_z + \frac{2\gamma}{\sqrt{N}}J_x(a^\dagger + a)
	\end{equation*}
	$J_{x,y,z} = \frac{1}{2}\sum_{k=1}^N \sigma_{x,y,z}^k$.
	\item Hamiltonian commutes with the parity operator $\Pi = \exP{i\pi\Lambda}$ with $\Lambda = a^\dagger a + J_z + j\mathds{1}$
	\item When $\gamma$ reaches the critical value $\gamma_c=\sqrt{\omega \omega_0}/2$, then there is a 2nd order phase transition from a normal phase $\gamma<\gamma_c$ with no photons in the ground state, to the superradiant phase $\gamma>\gamma_c$ where the ground state has a finite amount of photons and excited atoms.
	\item I observed something similar when computing the exact eigenstates of the Fransson model. Maybe it is related to the Dicke model, and maybe there is also a superradiant phase transition.
	\item However, there are important differences: Dicke model has one bosonic mode -  Franssons has $N-1$, the electron-boson couplings are very different, the electrons can hop in Franssons model
\end{itemize}
\section{Quantum scarring as a possible reason for not finding a steady state}
\begin{itemize}
	\item For dynamical systems in the chaotic regime, periodic orbits (=orbits that repeat themselves in time) are unstable
	\item Presence of unstable orbits is due to quantum scarring
	\item Quantum scarring: Restricts the degree of delocalization and leads to revivals in the dynamics
	\item 
\end{itemize}
\section{Ideas for annual convent talk}
\begin{itemize}
	\item Present model with sketches
	\item Show how Fransson's SOC potential looks like in real-space
	\item Show PES for the smaller molecules
\end{itemize}
\section{Model for CISS? How to reconcile transmission spectroscopy models with transport models}
\begin{itemize}
	\item Check out literature on photoelectron spectroscopy (theory): What are the common models used there?
	\item How do these models differ from the ones I use?
	\item Role of triplet states?
	\item Initial idea: Hamiltonian is of the form with three electronic states and two modes
	\begin{align*}
		H &= H_\mathrm{mol} + H_\mathrm{electrode} + H_\mathrm{mol-electrode} + H_\mathrm{vib} + H_\mathrm{el-vib}\\
		H_\mathrm{mol} &= \sum_{s=\{S_0,S_1,T\}} \varepsilon_s \vert s\rangle \langle s\vert + \lambda_0 (\vert S_1\rangle \langle T\vert +\vert T\rangle \langle S_1\vert) \\
		H_\mathrm{electrode} &= \sum_{k\in\{L,R\}} \varepsilon_k c_k^\dagger c_k^\pdagger \\
		H_\mathrm{mol-electrode} &= \sum_{k\in\{L,R\}} (V_k \vert k\rangle \langle S_1\vert + \mathrm{h.c.})\\
		H_\mathrm{vib} &= \sum_\alpha \omega_\alpha b_\alpha^\dagger b_\alpha^\pdagger\\
		H_\mathrm{el-vib} &=g (b_1^\dagger+ b_1^\pdagger)(\vert S_0\rangle\langle T\vert + \mathrm{h.c.})+  \lambda_1 (b_2^\dagger+ b_2^\pdagger)(\vert S_1\rangle\langle T\vert + \mathrm{h.c.})
	\end{align*}
	\item Strong light-matter coupling can enhance singlet fission ("Not dark yet for strong light-matter coupling to accelerate singlet fission dynamics") $\rightarrow$ similarly possible for CISS?
	\item Relationship between intersystem crossing and CISS?
\end{itemize}
\section{Stability of fixed points}
\begin{itemize}
	\item $L=2$ Peierls model
	\begin{itemize}
		\item Isolated
		\item Open
	\end{itemize}
\end{itemize}
\printbibliography
\end{document}


